% --- Archon physics formalization source begin ---
% archon:physics
% archon:chemistry
% archon:covers IChO2026Problems/problem_icho_2026_t6_a7.lean
% archon:source-report reports/icho_2026/problem_icho_2026_t6_a7.source.json
% archon:problem-id icho_2026_t6
% archon:part-id T6-A7
% archon:previous-part icho_2026_t6_a6
% archon:previous-part-policy natural_language_prerequisite_only

\chapter{Icho Chemistry problem icho\_2026\_t6\_a7}
\label{ch:IChO2026Problems_problem_icho_2026_t6_a7}

\paragraph{Problem source.}
\#\# Source
58th International Chemistry Olympiad, Tashkent, Uzbekistan, 2026.
Official-materials index: https://www.icho2026.uz/problems
English problem PDF: ../icho\_2026\_source/raw/theory\_problem.pdf
English solution/rubric PDF: ../icho\_2026\_source/raw/theory\_solution.pdf
Problem PDF page 56; printed local question page 5; solution starts on PDF page 58.
The page PNG is primary visual evidence for formulas, structures, tables, and figures.

\#\# T6. Carbon Nanorings
T6. Carbon Nanorings 7\% In 2019, C18, the first example of a new class of carbon allotropes called cyclo[n]carbons (C𝑛) was detected. Cyclocarbons are monocyclic, all-carbon compounds with high strain energy and low stability. They can be generated and analysed by bond-resolved Atomic Force Microscopy (AFM). Cyclocarbons have interesting aromatic properties based on their π systems. C13 has triplet (3C13, spin S = 1), and singlet (1C13, spin S = 0) forms. 6.1 Based on Hückel’s rule, fill in the table with the number of distinct aromatic (A) and anti-aromatic (AA) 𝜋-systems present in each of C18, C16, 3C13, and 1C13. Fill in all blanks. Later, cyclo[n]carbons C10, C12, C14, C20, and C26 were generated by a combination of dehalogenation and retro-Bergman (1) reactions on a surface. AFM images taken during the synthesis of C14 are shown below: 6.2 Draw structures B-D. Note they may contain one or more unpaired electrons. The structure of A is given as an example. 6.3 The voltages applied during the AFM-mediated synthesis of C𝑛may vary. Using the given C–X bond energies, choose the possible halogen(s) in the halogenated reagent for C18 synthesis via a retro-Bergman reaction with a voltage of 2.5 V acting on the electrons. Bond (C–X) Bond Dissociation Energy / kJ mol−1 C–F 467 C–Cl 346 C–Br 290 C–I 228 In 2025, the first relatively stable cyclo[48]carbon was synthesised. The molecule was stabilised by catenation (forming interlocking rings) with macrocycle E and was first characterised by electrospray mass spectrometry in positive mode. 6.4 In the mass spectrum obtained, four intense peaks were observed m/z: 591, 783, 879, and 1174. Suggest the identity of these ions. Use integer atomic masses and assume no fragmentation happened. The ion corresponding to m/z = 591 is given as an example. Cycloparaphenylene (CPP) is another cyclic, highly-strained molecule commonly called a carbon nanohoop. The synthesis of the smallest known member, [5]CPP, is shown below, where PhI(OAc)2 acts as a two electron oxidant. (equiv. = Equivalent, excess = Excess) 6.5 Draw the structures of F-L. Next generation nanobelts were constructed from porphyrin rings and triple bonds using the template synthesis method. Hint: M is the thermodynamic product and has four types of protons. (equiv. = Equivalent) 6.6 Draw the structures of M-R. When aromatic rings are joined to form a macrocycle, oxidation or reduction can produce global aromaticity, in which π electrons are delocalised around the entire macrocycle. Only the π system forming the continuous conjugated pathway around the ring participates in this delocalisation. As an example, [n]CPPs can become globally aromatic upon addition or removal of electrons. The number of π electrons is counted only along the bold pathway shown below. When P6 is oxidised, it can exhibit global aromaticity and global anti-aromaticity.

\#\# Current subquestion 6.7
Based on Hückel’s rule, find the minimum number of electrons, n(e), from P6 that should be removed to achieve global aromaticity. Write the total number of π electrons, n(t), in the global aromatic system.

\#\# Dataset metadata
IChO 2026; T6; raw rubric points=4; kind=theory; formalization\_ready=true.

\paragraph{Shared problem context.}
T6. Carbon Nanorings 7\% In 2019, C18, the first example of a new class of carbon allotropes called cyclo[n]carbons (C𝑛) was detected. Cyclocarbons are monocyclic, all-carbon compounds with high strain energy and low stability. They can be generated and analysed by bond-resolved Atomic Force Microscopy (AFM). Cyclocarbons have interesting aromatic properties based on their π systems. C13 has triplet (3C13, spin S = 1), and singlet (1C13, spin S = 0) forms. 6.1 Based on Hückel’s rule, fill in the table with the number of distinct aromatic (A) and anti-aromatic (AA) 𝜋-systems present in each of C18, C16, 3C13, and 1C13. Fill in all blanks. Later, cyclo[n]carbons C10, C12, C14, C20, and C26 were generated by a combination of dehalogenation and retro-Bergman (1) reactions on a surface. AFM images taken during the synthesis of C14 are shown below: 6.2 Draw structures B-D. Note they may contain one or more unpaired electrons. The structure of A is given as an example. 6.3 The voltages applied during the AFM-mediated synthesis of C𝑛may vary. Using the given C–X bond energies, choose the possible halogen(s) in the halogenated reagent for C18 synthesis via a retro-Bergman reaction with a voltage of 2.5 V acting on the electrons. Bond (C–X) Bond Dissociation Energy / kJ mol−1 C–F 467 C–Cl 346 C–Br 290 C–I 228 In 2025, the first relatively stable cyclo[48]carbon was synthesised. The molecule was stabilised by catenation (forming interlocking rings) with macrocycle E and was first characterised by electrospray mass spectrometry in positive mode. 6.4 In the mass spectrum obtained, four intense peaks were observed m/z: 591, 783, 879, and 1174. Suggest the identity of these ions. Use integer atomic masses and assume no fragmentation happened. The ion corresponding to m/z = 591 is given as an example. Cycloparaphenylene (CPP) is another cyclic, highly-strained molecule commonly called a carbon nanohoop. The synthesis of the smallest known member, [5]CPP, is shown below, where PhI(OAc)2 acts as a two electron oxidant. (equiv. = Equivalent, excess = Excess) 6.5 Draw the structures of F-L. Next generation nanobelts were constructed from porphyrin rings and triple bonds using the template synthesis method. Hint: M is the thermodynamic product and has four types of protons. (equiv. = Equivalent) 6.6 Draw the structures of M-R. When aromatic rings are joined to form a macrocycle, oxidation or reduction can produce global aromaticity, in which π electrons are delocalised around the entire macrocycle. Only the π system forming the continuous conjugated pathway around the ring participates in this delocalisation. As an example, [n]CPPs can become globally aromatic upon addition or removal of electrons. The number of π electrons is counted only along the bold pathway shown below. When P6 is oxidised, it can exhibit global aromaticity and global anti-aromaticity.

\paragraph{Current subquestion.}
Based on Hückel’s rule, find the minimum number of electrons, n(e), from P6 that should be removed to achieve global aromaticity. Write the total number of π electrons, n(t), in the global aromatic system.

\paragraph{Recorded answer/context.}
We should count electrons around the abbreviated structure of P6 as shown below: The total number of 𝜋-electrons is 14 × 6 = 84 in the ground state. To achieve global aromaticity by Hückel's rule, we need 4𝑛+ 2 𝜋-electrons. The minimum number of 𝜋-electrons that need to be lost is two electrons and the total number of 𝜋-electrons is 82. n(e) 2 n(t) 82 4 pts for fully correct. 2 pts if total number of electrons is wrong but it satisfies Hückel's 4𝑛+ 2 rule. No partial points for the loss of two electrons. 4.0 pt

\paragraph{Figure/image paths.}
\begin{itemize}
\item ../icho\_2026\_source/image/T6\_page-5.png
\item ../icho\_2026\_source/image/T6\_page-4.png
\end{itemize}

\paragraph{Reusable previous-part conclusions.}
\begin{itemize}
\item Source T6-A6. Question: Draw the structures of M-R. Reusable conclusions: Visual-answer notice: the official solution contains essential ticks, structures, tables, graphs, or equations that PDF plain-text extraction may omit. Consult the cited solution PDF page.

M N O Q R 20.0 pt RUBRICS: M: 4 pts Substitution pattern should be determined with information on four types of proton. 1 pt for any structure with one t-Bu group. 3 pt for structure with t-Bu added at both ortho positions, 1 pt for any other structure with two t-Bu groups. N: 4 pts 4 pts if amine is given instead of the aldehyde (Delepine reaction). Does not affect the grading of the rest of the parts. 1 pt if Br instead of aldehyde. O: 4 pts Q: 4 pts R: 4 pts - 1 pt if double bonds in porphyrin drawn incorrectly, this penalty will be applied at most once. - 2 pt if Zn is missing, this penalty will be applied at most once. Policy: natural\_language\_prerequisite\_only; the current target and later answers must not be assumed
\end{itemize}

\paragraph{Verified electron-count model.}
The highlighted pathway is represented by seven explicit two-electron
$\pi$ features per subunit.  Their sum is proved to be $14$; for six
subunits this gives $84$, and the least oxidation reaching a $4k+2$ count
removes two electrons, leaving $82$.

\paragraph{Formalization target.}
The verified Lean formalization is in `IChO2026Problems/problem\_icho\_2026\_t6\_a7.lean`,
with its theorem and lemma proofs complete.
The Lean declarations must preserve every mathematical object, quantifier, hypothesis, side condition, approximation/error guarantee, and requested conclusion from the source statement.
The implementation reuses the pinned Mathlib, Physlib, and Chemistry libraries,
with local definitions only for source-specific objects.

\begin{theorem}[Icho Chemistry formalization target]
\label{thm:physics:icho_2026_t6_a7:target}
\lean{Icho2026T6A7.icho_2026_t6_a7_target}
\leanok
The mapped declarations encode the stated calculation and its verified conclusion.
\end{theorem}
\begin{proof}
The proof unfolds the source-specific definitions and establishes the mapped
conclusion from the encoded hypotheses.
\end{proof}
% --- Archon physics formalization source end ---
